\chapter{进制与位运算}

\section{十进制与二进制}

人类最常用的计数方式是\textbf{十进制}，而计算机底层是\textbf{二进制}。

\subsection{简介}

十进制逢十进一，也就是说，在十进制中每当出现十就往前进一位，例如 $9 + 2 = 11$。在这个过程中，$9$ 先加上 $1$ 得到了十，要向十位进一位，得到 $10$，再加上 $1$，得到 $11$。

和十进制类似，二进制逢二进一，也就是说，在二进制中每当出现二就往前进一位，例如在二进制中，$1 + 1 = 10$。在这个过程中，$1$ 加上 $1$ 得到了二，要向第二位进一位，得到 $10$。

\subsection{二进制的运算规则}

二进制最基本的运算规则是逢二进一。在此基础上，有如下等式：
\[
\begin{cases}
1 + 1 &= 10\\
1 - 1 &= 0\\
1 \times 1 &= 1\\
1 \div 1 &= 1
\end{cases}
\]
基于这些等式，可以进行二进制的运算。

\subsection{二进制转十进制}

二进制转十进制的方法是：

对于长度为 $l$ 的二进制数从右往左每一位，从 $0$ 开始计数，将第 $i$ 位（$i$ 从 $0$ 开始计数）上的数字乘上 $2^i$ 并相加，即得到二进制数。那么若第 $i$ 位上的数值为 $s_i$，则十进制数 $x$ 的值为：
\[
x = \sum\limits_{i=0}^{l-1} s_i \times 2^i
\]

代码实现如下：

\begin{lstlisting}[caption=二进制转十进制]
#include <iostream>
#include <string>
#include <algorithm>
#include <cmath>
using namespace std;

string s;
int res;

int main() {
	
	cin >> s;
	
	reverse(s.begin(), s.end());
	
	for (int i=0; i<s.size(); i++) {
		res += (s[i] - '0') * pow(2, i);
	}
	
	cout << res;
	
	return 0;
}
\end{lstlisting}

\subsection{十进制转二进制}

对于一个十进制数，要将其转为二进制，通常采用\textbf{短除法}。具体的说，对于一个十进制数 $x$，将其不断地除以 $2$，直到除得的商为 $0$，将得到的余数倒序排列，即是得到的二进制数。

代码实现如下：

\begin{lstlisting}[caption=十进制转二进制]
#include <iostream>
#include <string>
using namespace std;

int x;
string res;

int main() {
	
	cin >> x;
	while (x) {
		res = (char)(x % 2 + '0') + res;
		x /= 2;
	}
	cout << res;
	
	return 0;
}
\end{lstlisting}

\subsection{C++ 中的二进制与十进制数}

在 C++ 中，可以通过 0b 在变量中存储二进制数，例如 0b1001 代表了二进制的 $1001$，即十进制 $9$。不加前缀，默认为十进制，例如 18 代表了十进制 $18$。

除此之外，C++ 提供了一种数据结构——bitset。下面通过一些代码来介绍 bitset 的使用：

\begin{lstlisting}[caption=bitset 的使用]
#include <bitset>
#include <string>
using namespace std;

int main() {
	
	// 创建一个 8 位的 bitset，初始化为 0
	bitset<8> b1;
	// 创建一个 8 位的 bitset，初始化为 10
	bitset<8> b2(10);
	// 创建一个 8 位的 bitset，初始化为二进制 00101010
	bitset<8> b3("00101010");
	
	return 0;
}
\end{lstlisting}

\section{八进制与十六进制}

除了二进制和十进制，在计算机世界中，八进制和十六进制也很常用。

\subsection{八进制的运算规则}

八进制逢八进一，也就是说，运算中出现八就要往前进一位。于是，有如下的等式：

\[
\begin{cases}
7 + 1 &= 10\\
10 - 1 &= 7\\
2 \times 4 &= 10\\
11 \div 3 &= 3
\end{cases}
\]

其中的 $10$ 对应十进制的 $8$，$11$ 对应十进制的 $9$。

\subsection{十六进制的表示和运算规则}

十六进制逢十六进一，数字的范围在零到十五之间。那么如何表示十及十以上的数字呢，我们使用字母 $a \sim f $ 表示，其中 $a$ 表示十进制 $10$，$b$ 表示十进制 $11$，一直到 $f$ 表示十进制 $15$。

基于以上规则，有如下等式：

\[
\begin{cases}
f + 1 &= 10\\
10 - 1 &= f\\
2 \times 8 &= 10\\
1b \div 3 &= 9
\end{cases}
\]

其中的 $f$ 对于十进制的 $15$，$10$ 对应十进制的 $16$，$1b$ 对于十进制的 $16 + 11 = 27$。

\subsection{C++ 中的八进制和十六进制数}

在 C++ 中，只需要以 0 开头，就可以表示八进制数。例如 01011 代表八进制的 $1011$（无符号），也就是十进制中的 $521$。只需要以 0x 开头，就可以表示十六进制，例如 0x3f 代表十六进制中的 $3f$（无符号），十进制中的 $63$。

在 C++ 中，0x3f3f3f3f，即十六进制 $3f3f3f3f$，是一个大于 $10^9$ 的数；0x7f7f7f7f，即十六进制 $7f7f7f7f$，是一个大于 $2 \times 10^9$ 的数。

\subsection{进制转换}

我们之前已经介绍了二进制与十进制的转换，其实，这些方法都是基本通用的。

对于其他进制转十进制，我们通常采用逐位相乘的方法。具体的说，对于 $n$ 进制转十进制，我们从右往左，从 $0$ 开始计数，将其第$i$ 位 $s_i$ 乘上 $n^i$ 相加，即可得到结果。那么若原 $n$ 进制数的长度为 $l$，则十进制数 $x$ 的值为：
\[
x = \sum_{i=0}^{l-1} s_i \times n^i
\]

对于十进制转其他进制，通常采用短除法。具体的说，对于十进制 $n$ 进制，将十进制数不断除以 $n$，直到商为 $0$，将得到的余数转化为 $n$ 进制中的一位数字后逆序排列即可。

由于各个进制对数的表示方法有差异，这里不给出统一的代码，读者可以针对具体需要，参考二进制与十进制转换的代码自行编写。

\section{计算机中数据的存储与表示}

\subsection{计算机中数据的存储}

计算机中的数据以二进制形式存储。在计算机中，一个二进制位称作一个比特（bit，简写 b，或称比特位，有时简称位）。每 8 个二进制位是一组，称为一个字节（byte，简写 B）。下面我们的介绍都以 1 个字节为例，实际编程中，一个 int 整型占用 4 个字节；一个 long long 长长整型占用 8 个字节。

计算机中的数分为有符号数和无符号数两种。有符号数从左往右第一个比特是符号为，0 代表正数。1 代表负数。无符号数从左往右第一个比特计入数字，默认正数。

\subsection{计算机中整型数据的表示}

计算机中数据的表示分为原码、反码、补码三种，计算机中的数据以补码形式存储。下面均以有符号八位整数举例。

\textbf{原码}是最便于人类阅读的一种表示方式。原码从左往右第一位是符号位，0 代表正数。1 代表负数。后面的是数字位，代表数的绝对值对应的二进制表示。例如 1000 0011 代表 $-3$。

\textbf{反码}是原码到补码转换的一个中间过程。正数的反码与原码相同；负数的反码在原码基础上，对除符号位以外的位取反。例如原码 1000 0011 对应的反码就是 1111 1100。

\textbf{补码}是计算机中数据的存储形式，最便于计算机进行数学计算。正数的补码与原码相同；负数的补码在反码基础上加 $1$。例如原码 1000 0011 对应的补码就是 1111 1101。

特别的，正零和负零的补码都是 0000 0000。

\section{位运算}

位运算是计算机处理数据的基本逻辑门。位运算包括按位与、按位或、按位非、按位异或（简作异或）、左移、右移。下面除右移外均以无符号八位整数举例。

\subsection{按位与、按位或}

\textbf{按位与}有两个运算数参与运算，其运算规则是，对两个运算数的每一位进行如下的运算：若给定的两个数都是 $1$，结果为 $1$；否则结果为 $0$。例如：二进制数 0110 0111 和 0001 0110 的按位与结果为 0000 0110。C++ 中符号为 \&。

\textbf{按位或}有两个运算数参与运算，其运算规则是，对两个运算数的每一位进行如下的运算：若给定的两个数有一个是 $1$，结果为 $1$；否则结果为 $0$。例如：二进制数 0110 0111 和 0001 0110 的按位或结果为 0111 1111。C++ 中符号为 |。

\subsection{按位非、异或}

\textbf{按位非}只有一个运算数参与运算，其运算规则是，对运算数的每一位进行如下的运算：若给定的数是 $1$，结果为 $0$；否则结果为 $1$。例如：二进制数 0110 0111 的按位非结果为 1001 1000。C++ 中符号为 $!$。

\textbf{按位异或}（简作\textbf{异或}）有两个运算数参与运算，其运算规则是，对两个运算数的每一位进行如下的运算：若给定的两个数相同，结果为 $0$，否则结果为 $1$。例如：二进制数 0110 0111 和 0001 0110 的按位或结果为 0111 0001。C++ 中符号为 \^{}，数学符号通常为 $\oplus$。

\subsection{左移、右移}

\textbf{左移}有两个运算数参与运算，用于将一个数的二进制向左移动若干位，其运算规则是：溢出的位数直接舍弃，空缺的位数用 0 填补。例如：二进制数 0110 0111 左移 2 位的结果为 1001 1100。C++ 中符号为 <<。

\textbf{右移}有两个运算数参与运算，用于将一个数的二进制向右移动若干位，其运算规则是：溢出的位数直接舍弃，空缺的位数用符号位填补（无符号数补 0）。例如：有符号二进制数 1110 0111 右移 2 位的结果为 1111 1001。C++ 中符号为 >>。